% Copyright 2025 by Robert Hildebrand
%This work is licensed under a
%Creative Commons Attribution-ShareAlike 4.0 International License (CC BY-SA 4.0)
%See http://creativecommons.org/licenses/by-sa/4.0/

\chapter{Metode Simpleks dalam Bentuk Tableau}
\label{ch:simplex-tableau}
\begin{outcome}
\begin{itemize}
\item Catat sebuah program linier dan bentuk slack-nya dalam \emph{tableau simpleks}
\item Baca solusi basis layak, nilai objektif, dan biaya tereduksi langsung dari tableau
\item Lakukan satu pivot simpleks sebagai rangkaian operasi baris elementer
\item Kenali syarat optimalitas dalam bentuk tableau dan baca solusi optimalnya
\end{itemize}
\end{outcome}

\begin{tryit}
\vizlink{simplex-tableau}{Pemutar Simpleks dalam Bentuk Tableau}: mekanisme pivot yang sama seperti dalam bab ini, dengan setiap operasi baris ditampilkan.

\vizlink{tableau-pivoter}{Pemutar Tableau Bebas}: sebuah arena Gauss--Jordan; lakukan pivot pada sembarang entri dan amati hasilnya.
\end{tryit}



% Fixed-width centred column for evenly spaced tableau entries.
% (Z is unused by array/tabularx/siunitx, so this is safe book-wide.)
\newcolumntype{Z}{>{\centering\arraybackslash}p{0.95cm}}

Metode kamus dari Bab~\ref{ch:simplex-dictionary} dan perhitungan matriks dari
Bab~\ref{ch:simplex-matrix} menjalankan aritmetika yang \emph{sama}; keduanya
hanya berbeda dalam seberapa banyak aritmetika itu kita tuliskan. \emph{Tableau
simpleks} adalah cara pencatatan yang paling ringkas: satu tabel bilangan yang
menyimpan bentuk slack program linier dan memperbaruinya langsung pada setiap
pivot. Semua yang diperlukan metode simpleks---variabel mana yang masuk,
variabel mana yang keluar, dan apakah proses telah selesai---dapat dibaca
langsung dari tabel tersebut.

Kami menggunakan program linier yang sama yang diselesaikan oleh kamus sebelumnya,
\begin{equation*}
\begin{aligned}
\max\quad & z = 2x + 3y\\
\st\quad  & x + y  \le 9,\\
          & 2x + y \le 16,\\
          & x + 2y \le 14,\\
          & x, y \ge 0,
\end{aligned}
\end{equation*}
sehingga Anda dapat membandingkan kedua penyajian itu baris demi baris.

\section{Membangun Tableau}
\label{sec:tableau-building}

Pertama, ubah program ke \emph{bentuk standar} dengan menambahkan variabel
slack tak negatif pada setiap kendala:
\begin{equation*}
\begin{aligned}
x + y  + s_1 &= 9,\\
2x + y + s_2 &= 16,\\
x + 2y + s_3 &= 14,
\end{aligned}
\qquad x,y,s_1,s_2,s_3 \ge 0.
\end{equation*}
Kita juga mencatat objektif sebagai persamaan,
\[
z - 2x - 3y = 0,
\]
agar objektif berada dalam tableau dengan kedudukan yang sama seperti kendala.
Persamaan ini adalah \emph{baris objektif} (atau baris $z$); kita menempatkannya
di bagian atas.

Setiap baris tableau memuat koefisien satu persamaan, dan kolom terakhir memuat
ruas kanannya. Kolom paling kiri menamai \emph{variabel basis} pada baris
tersebut: variabel yang muncul dengan koefisien $1$ pada baris itu dan $0$ pada
setiap baris lainnya. Dari tableau awal, kita membaca
\begin{center}
\renewcommand{\arraystretch}{1.3}
\begin{tabular}{c|ZZZZZ|Z}
\toprule
Basis & $x$ & $y$ & $s_1$ & $s_2$ & $s_3$ & RHS\\
\midrule
$z$   & $-2$ & $-3$ & $0$ & $0$ & $0$ & $0$\\
\midrule
$s_1$ & $1$ & $1$ & $1$ & $0$ & $0$ & $9$\\
$s_2$ & $2$ & $1$ & $0$ & $1$ & $0$ & $16$\\
$s_3$ & $1$ & $2$ & $0$ & $0$ & $1$ & $14$\\
\bottomrule
\end{tabular}
\end{center}
bahwa basisnya adalah $\{s_1,s_2,s_3\}$. Menetapkan variabel nonbasis
$x=y=0$ memberikan solusi basis layak $s_1=9,\ s_2=16,\ s_3=14$, yang
dibaca langsung dari kolom RHS, dengan nilai objektif $z=0$ (RHS pada baris
$z$). Entri baris $z$, yaitu $-2$ dan $-3$, mencatat \emph{biaya tereduksi}
$x$ dan $y$ dengan tanda yang dibalik: menaikkan $x$ atau $y$ dari nol akan
mengubah $z$ masing-masing dengan laju $+2$ dan $+3$, yakni negatif dari entri
yang disimpan.

\begin{steppanel}{algpivot}
\label{panel:tableau-sign-convention}%
\slab{algpivot}{114}{Konvensi tanda} Dalam bentuk kamus, kita menuliskan
$z = 2x + 3y$ dan menaikkan variabel dengan koefisien \emph{positif terbesar}.
Tableau justru menyimpan baris $z - 2x - 3y = 0$, sehingga koefisien yang sama
muncul dengan \emph{tanda berlawanan}. Karena itu, kedua aturan tampak seperti
cerminan dan menghasilkan keputusan yang identik: dalam bentuk tableau, kita
memasukkan variabel dengan entri baris $z$ yang \emph{paling negatif}, dan kita
berhenti ketika \emph{setiap} entri baris $z$ tak negatif.
\end{steppanel}

\section{Satu Pivot, Tiga Warna}
\label{sec:tableau-pivot}

Sebuah pivot memiliki tiga langkah, dan kita mewarnainya secara konsisten di
seluruh buku: {\color{algenter}\textbf{hijau}} memilih variabel masuk,
{\color{algleave}\textbf{oranye}} menjalankan uji rasio untuk menemukan
variabel keluar, dan {\color{algpivot}\textbf{biru}} menandai elemen pivot di
tempat keduanya berpotongan.

\begin{center}
\renewcommand{\arraystretch}{1.3}
\begin{tabular}{c|ZZZZZ|Z|c}
\toprule
Basis & $x$ & \cellcolor{algenter!18}$y$ & $s_1$ & $s_2$ & $s_3$ & RHS & Rasio\\
\midrule
$z$   & $-2$ & \cellcolor{algenter!18}$-3$ & $0$ & $0$ & $0$ & $0$ & \\
\midrule
$s_1$ & $1$ & \cellcolor{algenter!18}$1$ & $1$ & $0$ & $0$ & $9$  & $9/1=9$\\
$s_2$ & $2$ & \cellcolor{algenter!18}$1$ & $0$ & $1$ & $0$ & $16$ & $16/1=16$\\
\rowcolor{algleave!15}$s_3$ & $1$ & \cellcolor{algpivot!30}$2$ & $0$ & $0$ & $1$ & $14$ & $\mathbf{14/2=7}$\ \ding{47}\\
\bottomrule
\end{tabular}
\end{center}

\noindent%
{\color{algenter}\ding{43}\ \textbf{Masuk:}} $y$ (entri baris $z$ sebesar $-3$, yang paling negatif).\quad
{\color{algleave}\ding{47}\ \textbf{Keluar:}} $s_3$ (perbandingan terkecil $7$).\quad
{\color{algpivot}\ding{114}\ \textbf{Pivot}} pada elemen biru.

\medskip
Pivot itu sendiri hanyalah eliminasi Gauss pada kolom pivot, yang mengubahnya
menjadi kolom satuan dengan $1$ pada baris pivot.

\begin{steppanel}{algenter}
\slab{algenter}{43}{Masuk} Entri baris $z$ adalah $-2$ (untuk $x$) dan $-3$
(untuk $y$). Entri yang paling negatif adalah $-3$, sehingga $y$
\textbf{masuk}: menaikkan $y$ memperbaiki $z$ paling cepat.
\end{steppanel}
\begin{steppanel}{algleave}
\slab{algleave}{47}{Uji rasio} Bagi RHS setiap kendala dengan entri
\emph{positif} pada kolom $y$:
\[
\frac{9}{1}=9,\qquad \frac{16}{1}=16,\qquad \frac{14}{2}=7\ (\text{terkecil}).
\]
Rasio terkecil, $7$, terjadi pada baris $s_3$, sehingga $s_3$ \textbf{keluar}.
Elemen pivot adalah entri tempat kolom masuk bertemu dengan baris keluar, yaitu
$2$.
\end{steppanel}
\begin{steppanel}{algpivot}
\slab{algpivot}{114}{Pivot} Skalakan baris pivot sehingga elemen pivot menjadi $1$, kemudian hapus kolom $y$ di tempat lain:
\[
\begin{aligned}
\text{baris }y\text{ baru} &\leftarrow \tfrac12\,(\text{baris }s_3) = [\,\tfrac12,\ 1,\ 0,\ 0,\ \tfrac12 \mid 7\,],\\
\text{baris }z &\leftarrow \text{baris }z + 3\,(\text{baris }y\text{ baru}),\\
\text{baris }s_1 &\leftarrow \text{baris }s_1 - 1\,(\text{baris }y\text{ baru}),\\
\text{baris }s_2 &\leftarrow \text{baris }s_2 - 1\,(\text{baris }y\text{ baru}).
\end{aligned}
\]
\end{steppanel}

\noindent Hasilnya adalah tableau kedua, dengan basis $\{s_1,s_2,y\}$:
\begin{center}
\renewcommand{\arraystretch}{1.3}
\begin{tabular}{c|ZZZZZ|Z|c}
\toprule
Basis & \cellcolor{algenter!18}$x$ & $y$ & $s_1$ & $s_2$ & $s_3$ & RHS & Rasio\\
\midrule
$z$   & \cellcolor{algenter!18}$-\tfrac12$ & $0$ & $0$ & $0$ & $\tfrac32$ & $21$ & \\
\midrule
\rowcolor{algleave!15}$s_1$ & \cellcolor{algpivot!30}$\tfrac12$ & $0$ & $1$ & $0$ & $-\tfrac12$ & $2$ & $\mathbf{2/(1/2)=4}$\ \ding{47}\\
$s_2$ & \cellcolor{algenter!18}$\tfrac32$ & $0$ & $0$ & $1$ & $-\tfrac12$ & $9$ & $9/(3/2)=6$\\
$y$   & \cellcolor{algenter!18}$\tfrac12$ & $1$ & $0$ & $0$ & $\tfrac12$ & $7$ & $7/(1/2)=14$\\
\bottomrule
\end{tabular}
\end{center}

\noindent%
{\color{algenter}\ding{43}\ \textbf{Masuk:}} $x$ (satu-satunya entri baris $z$ yang negatif, $-\tfrac12$).\quad
{\color{algleave}\ding{47}\ \textbf{Keluar:}} $s_1$ (perbandingan terkecil $4$).\quad
{\color{algpivot}\ding{114}\ \textbf{Pivot}} pada $\tfrac12$.

\medskip
Perhatikan bahwa nilai objektif telah naik dari $0$ menjadi $21$. Satu pivot
lagi pada elemen biru (mengalikan baris $s_1$ dengan $2$, lalu menolkan kolom
$x$) menghasilkan tableau ketiga, dengan basis $\{x,s_2,y\}$:
\begin{center}
\renewcommand{\arraystretch}{1.3}
\begin{tabular}{c|ZZZZZ|Z}
\toprule
Basis & $x$ & $y$ & $s_1$ & $s_2$ & $s_3$ & RHS\\
\midrule
\rowcolor{algopt!12}$z$ & $0$ & $0$ & $1$ & $0$ & $1$ & $23$\\
\midrule
$x$   & $1$ & $0$ & $2$  & $0$ & $-1$ & $4$\\
$s_2$ & $0$ & $0$ & $-3$ & $1$ & $1$  & $3$\\
$y$   & $0$ & $1$ & $-1$ & $0$ & $1$  & $5$\\
\bottomrule
\end{tabular}
\end{center}

\section{Optimalitas dan Membaca Solusi}
\label{sec:tableau-optimality}

Setiap entri baris $z$ sekarang tak negatif, sehingga tidak ada variabel yang
dapat masuk dan memperbaiki objektif: tableau tersebut \textbf{optimal}. Kita
membaca solusi langsung dari kolom RHS dengan menetapkan setiap variabel basis
sama dengan ruas kanannya dan setiap variabel nonbasis sama dengan nol:
\[
\boxed{\,x = 4,\quad y = 5,\quad z = 23\,}
\qquad\text{dengan } s_2 = 3,\ s_1 = s_3 = 0.
\]
Hasil ini tepat sama dengan solusi kamus, sebagaimana mestinya. Slack tak nol
$s_2=3$ menunjukkan bahwa kendala kedua $2x+y\le 16$ memiliki kelonggaran
(memang $2(4)+5 = 13 < 16$), sedangkan $s_1=s_3=0$ menunjukkan bahwa kendala
pertama dan ketiga aktif pada optimum.

\begin{algocard}{Algoritma: Metode Simpleks (Bentuk Tableau)}
\algostep{algenter}{\ding{43}}{Pilih variabel masuk}{Periksa baris objektif. Pilih variabel nonbasis dengan entri \emph{negatif}; pilihan umum adalah entri yang paling negatif. Jika tidak ada entri negatif, berhenti: tableau saat ini optimal.}
\algostep{algleave}{\ding{47}}{Uji rasio untuk variabel keluar}{Untuk setiap baris dengan entri \emph{positif} pada kolom masuk, bentuk rasio $\dfrac{\text{ruas kanan}}{\text{entri kolom masuk}}$. Baris dengan rasio terkecil keluar. (Jika tidak ada entri positif, masalah tersebut tak terbatas.)}
\algostep{algpivot}{\ding{114}}{Pivot}{Kalikan baris pivot agar elemen pivot menjadi $1$, lalu tambahkan kelipatan baris itu ke setiap baris lain, termasuk baris objektif, agar entri lain pada kolom masuk menjadi $0$.}
\algostep{algopt}{\ding{52}}{Ulangi hingga optimal}{Ulangi sampai setiap entri baris objektif tak negatif. Kemudian baca solusi optimal dan nilai objektif dari kolom RHS.}
\end{algocard}

Tableau, kamus, dan bentuk matriks adalah tiga jendela menuju algoritma yang
sama: tableau paling cepat untuk perhitungan tangan, kamus menyatakan persamaan
yang mendasarinya secara eksplisit, dan bentuk matriks
(Bab~\ref{ch:simplex-matrix}) menunjukkan mengapa setiap pembaruan merupakan
perubahan basis.

\section{Variabel Artifisial dan Metode Big-M}
\label{sec:tableau-bigm}

Setiap tableau sejauh ini berawal dari basis slack, tempat variabel slack
menyediakan blok identitas siap pakai dan titik awal layak. Kendala ``$\ge$''
merusak kemudahan itu: variabel surplusnya masuk dengan koefisien $-1$, sehingga
tidak dapat berperan sebagai variabel basis, dan titik asal tidak layak. Bentuk
tableau menangani hal ini tepat seperti bentuk kamus
(Bab~\ref{ch:simplex-dictionary}), yaitu dengan menambahkan \emph{variabel
artifisial} pada setiap kendala semacam itu dan memberinya penalti berupa
konstanta besar $M$ dalam objektif (\emph{metode Big-M}). Satu-satunya hal baru
adalah pencatatannya: karena setiap variabel artifisial bermula \emph{di dalam}
basis, kita harus terlebih dahulu menolkan kolomnya pada baris objektif sebelum
aturan simpleks dapat diterapkan.

\subsection{Masalah tanpa Basis Slack Awal}

Kembali ke program linier
\[
\begin{aligned}
\max\quad & 2x + 3y \\
\st\quad & 2x + y \geq 5,\quad 2x + y \leq 16,\quad x + 2y \leq 14,\quad x,y\ge 0.
\end{aligned}
\]
Kendala pertama memerlukan variabel surplus $e_1$ dan variabel artifisial
$a_1$; kedua kendala lainnya menggunakan variabel slack biasa:
\[
\begin{aligned}
2x + y - e_1 + a_1 &= 5,\\
2x + y + s_2 &= 16,\\
x + 2y + s_3 &= 14,
\end{aligned}
\qquad
\max\ z = 2x + 3y - M a_1 .
\]

\begin{steppanel}{algslate}
\cstep{1}{Nolkan variabel artifisial pada baris objektif} Menuliskan objektif
sebagai $z - 2x - 3y + M a_1 = 0$ akan meninggalkan entri tak nol di bawah
variabel \emph{basis} $a_1$. Mengurangkan $M$ kali baris $a_1$ menghasilkan
tableau yang benar, dengan konsekuensi munculnya $M$ pada entri-entri lain:
\end{steppanel}
\begin{center}
\renewcommand{\arraystretch}{1.3}\small\setlength{\tabcolsep}{4.5pt}
\begin{tabular}{c|cccccc|c|c}
\toprule
Basis & $x$ & $y$ & $e_1$ & $s_2$ & $s_3$ & $a_1$ & RHS & Rasio\\
\midrule
$z$   & \cellcolor{algenter!18}$-2-2M$ & $-3-M$ & $M$ & $0$ & $0$ & $0$ & $-5M$ & \\
\midrule
\rowcolor{algleave!15}$a_1$ & \cellcolor{algpivot!30}$2$ & $1$ & $-1$ & $0$ & $0$ & $1$ & $5$  & $\mathbf{5/2=2.5}$\ \ding{47}\\
$s_2$ & \cellcolor{algenter!18}$2$ & $1$ & $0$ & $1$ & $0$ & $0$ & $16$ & $16/2=8$\\
$s_3$ & \cellcolor{algenter!18}$1$ & $2$ & $0$ & $0$ & $1$ & $0$ & $14$ & $14/1=14$\\
\bottomrule
\end{tabular}
\end{center}

\noindent%
{\color{algenter}\ding{43}\ \textbf{Masuk:}} $x$ (entri baris $z$, $-2-2M$, adalah yang paling negatif untuk $M$ besar).\quad
{\color{algleave}\ding{47}\ \textbf{Keluar:}} $a_1$ (perbandingan terkecil $2.5$).\quad
{\color{algpivot}\ding{114}\ \textbf{Pivot}} pada $2$.

\medskip
Pivot pertama mendorong variabel artifisial keluar dari basis. Setelah
melakukan pivot pada elemen biru, kita memperoleh
\begin{center}
\renewcommand{\arraystretch}{1.3}\small\setlength{\tabcolsep}{4.5pt}
\begin{tabular}{c|cccccc|c}
\toprule
Basis & $x$ & $y$ & $e_1$ & $s_2$ & $s_3$ & $a_1$ & RHS\\
\midrule
$z$   & $0$ & $-2$ & $-1$ & $0$ & $0$ & $1+M$ & $5$\\
\midrule
$x$   & $1$ & $\tfrac12$ & $-\tfrac12$ & $0$ & $0$ & $\tfrac12$ & $\tfrac52$\\
$s_2$ & $0$ & $0$ & $1$ & $1$ & $0$ & $-1$ & $11$\\
$s_3$ & $0$ & $\tfrac32$ & $\tfrac12$ & $0$ & $1$ & $-\tfrac12$ & $\tfrac{23}{2}$\\
\bottomrule
\end{tabular}
\end{center}
Variabel artifisial $a_1$ sekarang nonbasis, dan entri baris objektifnya $1+M$
positif, sehingga variabel itu tidak akan pernah masuk kembali: kita dapat
menghapus kolom $a_1$ dan melupakannya. Yang tersisa adalah tableau layak biasa,
yang kita bawa menuju optimalitas dengan aturan simpleks dari bagian sebelumnya.
Beberapa pivot kemudian, setiap entri baris objektif menjadi tak negatif:
\begin{center}
\renewcommand{\arraystretch}{1.3}\small\setlength{\tabcolsep}{4.5pt}
\begin{tabular}{c|ccccc|c}
\toprule
Basis & $x$ & $y$ & $e_1$ & $s_2$ & $s_3$ & RHS\\
\midrule
\rowcolor{algopt!12}$z$ & $0$ & $0$ & $0$ & $\tfrac13$ & $\tfrac43$ & $24$\\
\midrule
$x$   & $1$ & $0$ & $0$ & $\tfrac23$  & $-\tfrac13$ & $6$\\
$e_1$ & $0$ & $0$ & $1$ & $1$         & $0$         & $11$\\
$y$   & $0$ & $1$ & $0$ & $-\tfrac13$ & $\tfrac23$  & $4$\\
\bottomrule
\end{tabular}
\end{center}
Membaca tableau optimal memberikan $x = 6$, $y = 4$, dan $z = 24$, dengan
surplus $e_1 = 11$ pada kendala pertama, sedangkan kedua kendala berikutnya
aktif ($s_2 = s_3 = 0$).

\subsection{Mendeteksi Ketaklayakan dalam Bentuk Tableau}

Metode Big-M juga menunjukkan ketika suatu masalah \emph{tidak memiliki}
solusi layak: variabel artifisial menolak keluar dari basis. Pertimbangkan
\[
\begin{aligned}
\max\quad & 2x + 3y \\
\st\quad & 2x + y \geq 22,\quad 2x + y \leq 16,\quad x + 2y \leq 14,\quad x,y\ge 0,
\end{aligned}
\]
yang dua kendala pertamanya mengharuskan $2x+y$ sekaligus sedikitnya $22$ dan
paling banyak $16$. Menyusunnya tepat seperti sebelumnya dan menolkan variabel
artifisial pada baris objektif memberikan tableau awal
\begin{center}
\renewcommand{\arraystretch}{1.3}\small\setlength{\tabcolsep}{4.5pt}
\begin{tabular}{c|cccccc|c}
\toprule
Basis & $x$ & $y$ & $e_1$ & $s_2$ & $s_3$ & $a_1$ & RHS\\
\midrule
$z$   & $-2-2M$ & $-3-M$ & $M$ & $0$ & $0$ & $0$ & $-22M$\\
\midrule
$a_1$ & $2$ & $1$ & $-1$ & $0$ & $0$ & $1$ & $22$\\
$s_2$ & $2$ & $1$ & $0$ & $1$ & $0$ & $0$ & $16$\\
$s_3$ & $1$ & $2$ & $0$ & $0$ & $1$ & $0$ & $14$\\
\bottomrule
\end{tabular}
\end{center}
Menjalankan metode simpleks sampai optimum pada masalah berpenalti (dua pivot
sudah cukup) menghasilkan
\begin{center}
\renewcommand{\arraystretch}{1.3}\small\setlength{\tabcolsep}{4.5pt}
\begin{tabular}{c|cccccc|c}
\toprule
Basis & $x$ & $y$ & $e_1$ & $s_2$ & $s_3$ & $a_1$ & RHS\\
\midrule
\rowcolor{algleave!15}$z$ & $0$ & $0$ & $M$ & $\tfrac13+M$ & $\tfrac43$ & $0$ & $24-6M$\\
\midrule
$a_1$ & $0$ & $0$ & $-1$ & $-1$ & $0$ & $1$ & $6$\\
$x$   & $1$ & $0$ & $0$ & $\tfrac23$ & $-\tfrac13$ & $0$ & $6$\\
$y$   & $0$ & $1$ & $0$ & $-\tfrac13$ & $\tfrac23$ & $0$ & $4$\\
\bottomrule
\end{tabular}
\end{center}
Setiap entri baris objektif tak negatif, sehingga tableau ini optimal untuk
masalah Big-M. Namun, variabel artifisial $a_1$ masih menjadi variabel basis
dengan nilai positif $a_1 = 6$. Karena solusi sejati masalah asli akan
memungkinkan setiap variabel artifisial ditetapkan sama dengan nol, variabel
artifisial positif dalam tableau Big-M yang optimal menandakan bahwa
\textbf{program linier asli tidak layak}. Suku $-6M$ pada nilai objektif
$24-6M$ menyatakan hal yang sama: tidak ada objektif berhingga yang dapat
menghindari penalti tersebut.

\begin{center}
\renewcommand{\arraystretch}{1.3}
\begin{tabular}{@{}p{0.46\textwidth}p{0.46\textwidth}@{}}
\toprule
\textbf{Variabel artifisial keluar dari basis} & \textbf{Variabel artifisial tetap dalam basis ($>0$)}\\
\midrule
Basis layak untuk masalah asli telah ditemukan; hapus kolom artifisial dan lanjutkan. & Masalah asli tak layak; tidak ada basis layak. \\
\bottomrule
\end{tabular}
\par\captionof{table}{Interpretasi tableau akhir Big-$M$: variabel artifisial keluar dari basis versus tetap dalam basis.}\label{tab:simplex-tableau-tab01}% tab-num-auto

\end{center}

\section{Latihan}

\exgroup{Pemanasan}

\begin{ex}{Membaca tableau di tengah proses penyelesaian}{}
\label{ex:tableau-read-mid}
\stars{1} Tableau di bawah ini adalah tableau kedua dari contoh berjalan dalam
bab ini, yang diperoleh setelah satu pivot:
\begin{center}
\renewcommand{\arraystretch}{1.3}
\begin{tabular}{c|ZZZZZ|Z}
\toprule
Basis & $x$ & $y$ & $s_1$ & $s_2$ & $s_3$ & RHS\\
\midrule
$z$   & $-\tfrac12$ & $0$ & $0$ & $0$ & $\tfrac32$ & $21$\\
\midrule
$s_1$ & $\tfrac12$ & $0$ & $1$ & $0$ & $-\tfrac12$ & $2$\\
$s_2$ & $\tfrac32$ & $0$ & $0$ & $1$ & $-\tfrac12$ & $9$\\
$y$   & $\tfrac12$ & $1$ & $0$ & $0$ & $\tfrac12$ & $7$\\
\bottomrule
\end{tabular}
\end{center}
\begin{enumerate}
\item Variabel mana yang basis dan mana yang nonbasis?
\item Baca solusi basis layak dan nilai objektif.
\item Apa biaya tereduksi dari variabel nonbasis?
\item Apakah tableau ini optimal? Jika tidak, variabel mana yang masuk berikutnya?
\end{enumerate}
\exrefs{\S\ref{sec:tableau-building}, tableau terurai pertama; \S\ref{sec:tableau-pivot}}
\end{ex}

\begin{ex}{Membaca tableau}{}
\label{ex:tableau-reading}
\stars{1} Tableau berikut muncul saat menyelesaikan masalah maksimisasi, dengan
baris objektif di bagian atas:
\begin{center}
\renewcommand{\arraystretch}{1.3}
\begin{tabular}{c|ZZZZ|Z}
\toprule
Basis & $x_1$ & $x_2$ & $s_1$ & $s_2$ & RHS\\
\midrule
$z$   & $0$ & $0$ & $3$ & $1$ & $40$\\
\midrule
$x_1$ & $1$ & $0$ & $2$  & $-1$ & $6$\\
$x_2$ & $0$ & $1$ & $-1$ & $1$  & $4$\\
\bottomrule
\end{tabular}
\end{center}
Identifikasi variabel basis dan nonbasis, nyatakan solusi basis layak yang
bersesuaian beserta nilai objektifnya, dan jelaskan bagaimana baris objektif
menunjukkan bahwa tableau ini optimal.
\exrefs{\S\ref{sec:tableau-building}; \S\ref{sec:tableau-optimality}}
\end{ex}

\begin{ex}{Variabel masuk, uji rasio, dan elemen pivot}{}
\label{ex:tableau-ratio-test}
\stars{1} Pertimbangkan tableau
\begin{center}
\renewcommand{\arraystretch}{1.3}
\begin{tabular}{c|ZZZZZ|Z}
\toprule
Basis & $x_1$ & $x_2$ & $s_1$ & $s_2$ & $s_3$ & RHS\\
\midrule
$z$   & $-4$ & $-3$ & $0$ & $0$ & $0$ & $0$\\
\midrule
$s_1$ & $2$ & $3$ & $1$ & $0$ & $0$ & $6$\\
$s_2$ & $4$ & $1$ & $0$ & $1$ & $0$ & $8$\\
$s_3$ & $3$ & $4$ & $0$ & $0$ & $1$ & $12$\\
\bottomrule
\end{tabular}
\end{center}
Tentukan variabel masuk, lakukan uji rasio dengan menampilkan setiap rasio,
sebutkan variabel keluar, dan lingkari elemen pivot.
\exrefs{\S\ref{sec:tableau-pivot}}
\end{ex}

\exgroup{Soal inti}

\begin{ex}{Melakukan satu pivot}{}
\label{ex:tableau-one-pivot}
\stars{2} Mulailah dari tableau pada Latihan~\ref{ex:tableau-ratio-test},
lakukan pivot yang telah Anda identifikasi, lalu tuliskan tableau hasilnya.
Periksa bahwa kolom masuk telah menjadi kolom satuan dan nilai objektif pada
baris $z$ telah meningkat.
\exrefs{\S\ref{sec:tableau-pivot}}
\end{ex}

\begin{ex}{Menyelesaikan program linier dalam bentuk tableau}{}
\label{ex:tableau-solve-lp}
\stars{2} Gunakan metode simpleks dalam bentuk tableau untuk menyelesaikan
\[
\begin{aligned}
\max\quad & z = x + 2y\\
\st\quad  & x + y  \le 4,\\
          & x + 3y \le 6,\\
          & x, y \ge 0.
\end{aligned}
\]
Tampilkan setiap tableau; beri warna atau label pada kolom masuk, baris keluar,
dan elemen pivot di setiap langkah. Nyatakan solusi optimal dan nilai
objektifnya. (Anda seharusnya mencapai optimum dalam dua pivot.)
\exrefs{\S\ref{sec:tableau-pivot}; \S\ref{sec:tableau-optimality}}
\end{ex}

\begin{ex}{Penyelesaian lengkap dengan dua pivot}{}
\label{ex:tableau-two-pivot-solve}
\stars{2} Selesaikan program linier berikut dalam bentuk tableau dengan
menggunakan aturan entri paling negatif untuk memilih variabel masuk:
\[
\begin{aligned}
\max\quad & z = 3x + 2y\\
\st\quad  & 2x + y  \le 10,\\
          & x + 3y \le 15,\\
          & x, y \ge 0.
\end{aligned}
\]
Tampilkan tableau awal dan setiap tableau setelah pivot, serta sertakan solusi
optimal, nilai objektif optimal, dan nilai variabel slack pada maksimum.
\exrefs{\S\ref{sec:tableau-pivot}; \S\ref{sec:tableau-optimality}}
\end{ex}

\begin{ex}{Menyusun tableau Big-M}{}
\label{ex:tableau-bigm-setup}
\stars{2} Pertimbangkan program linier
\[
\begin{aligned}
\max\quad & 4x + 3y \\
\st\quad  & x + y \ge 4,\quad x + 2y \le 10,\quad x,y \ge 0.
\end{aligned}
\]
Tambahkan variabel surplus $e_1$ dan variabel artifisial $a_1$ untuk kendala
pertama, serta variabel slack $s_2$ untuk kendala kedua; lalu tuliskan objektif
sebagai $\max\ 4x + 3y - M a_1$. Susun tableau awal dengan baris objektif di
bagian atas, kemudian \emph{nolkan} variabel artifisial $a_1$ pada baris
objektif agar entrinya menjadi nol. Variabel mana yang masuk terlebih dahulu,
dan apakah variabel artifisial itu keluar pada pivot pertama?
\exrefs{\S\ref{sec:tableau-bigm}}
\end{ex}

\exgroup{Konsep dan keterkaitan}

\begin{ex}{Konvensi tanda pada baris $z$}{}
\label{ex:tableau-sign-convention}
\stars{2} Bentuk kamus pada Bab~\ref{ch:simplex-dictionary} memilih variabel
masuk dengan koefisien objektif \emph{positif}, sedangkan tableau memilih
variabel dengan entri baris $z$ yang \emph{negatif}.
\begin{enumerate}
\item Jelaskan asal pembalikan tanda dengan menggunakan fakta bahwa tableau menyimpan objektif sebagai persamaan $z - \c^\top\x = 0$.
\item Sebuah kamus berbunyi $z = 12 + 3x_1 - 2x_2$. Tuliskan baris $z$ tableau yang bersesuaian, termasuk entri RHS-nya.
\item Jelaskan mengapa entri RHS pada baris $z$ sama dengan nilai objektif saat ini, meskipun tanda setiap koefisien lainnya dibalik.
\end{enumerate}
\exrefs{\S\ref{sec:tableau-building}, panel konvensi tanda}
\end{ex}

\begin{ex}{Dari tableau ke kamus}{}
\label{ex:tableau-to-dictionary}
\stars{2} Tuliskan kamus yang bersesuaian dengan tableau pada
Latihan~\ref{ex:tableau-reading}, dengan menyatakan $z$ dan setiap variabel
basis dalam variabel nonbasis $s_1$ dan $s_2$. Pastikan bahwa menetapkan
variabel nonbasis sama dengan nol mereproduksi solusi basis layak yang sama.
\exrefs{\S\ref{sec:tableau-building}; Bab~\ref{ch:simplex-dictionary}}
\end{ex}

\begin{ex}{Mengenali masalah tak terbatas}{}
\label{ex:tableau-unbounded}
\stars{2} Saat memecahkan masalah maksimasi Anda mendapatkan
\begin{center}
\renewcommand{\arraystretch}{1.3}
\begin{tabular}{c|ZZZZ|Z}
\toprule
Basis & $x_1$ & $x_2$ & $s_1$ & $s_2$ & RHS\\
\midrule
$z$   & $-2$ & $0$ & $0$ & $1$ & $10$\\
\midrule
$x_2$ & $-1$ & $1$ & $0$ & $1$ & $3$\\
$s_1$ & $-3$ & $0$ & $1$ & $2$ & $5$\\
\bottomrule
\end{tabular}
\end{center}
Variabel mana yang harus masuk? Cobalah uji rasio dan jelaskan mengapa uji itu
gagal. Apa yang ditunjukkan kegagalan ini tentang program linier tersebut?
\exrefs{\S\ref{sec:tableau-optimality}, langkah uji rasio pada kartu algoritma}
\end{ex}

\begin{ex}{Membaca ketaklayakan}{}
\label{ex:tableau-infeasible}
\stars{2} Misalkan metode Big-M yang diterapkan pada suatu masalah maksimisasi
berakhir (setiap entri baris objektif tak negatif) dengan tableau
\begin{center}
\renewcommand{\arraystretch}{1.3}\small\setlength{\tabcolsep}{4.5pt}
\begin{tabular}{c|ccccc|c}
\toprule
Basis & $x$ & $y$ & $e_1$ & $s_2$ & $a_1$ & RHS\\
\midrule
$z$   & $0$ & $0$ & $M$ & $2$ & $0$ & $8-3M$\\
\midrule
$a_1$ & $0$ & $0$ & $-1$ & $-1$ & $1$ & $3$\\
$x$   & $1$ & $0$ & $0$ & $1$ & $0$ & $5$\\
$y$   & $0$ & $1$ & $-1$ & $0$ & $0$ & $2$\\
\bottomrule
\end{tabular}
\end{center}
Apakah tableau tersebut optimal untuk masalah berpenalti? Berapakah nilai
variabel artifisial $a_1$, dan apa yang dinyatakannya tentang program linier
asli? Jelaskan peran suku $-3M$ dalam nilai objektif.
\exrefs{\S\ref{sec:tableau-bigm}}
\end{ex}

\exgroup{Soal tantangan}

\begin{ex}{Menyertifikasi ketakterbatasan dari tableau}{}
\label{ex:tableau-unbounded-certificate}
\stars{3} Menggunakan metode simpleks pada
\[
\begin{aligned}
\max\quad & z = x_1 + x_2\\
\st\quad  & x_1 - x_2  \le 1,\\
          & -x_1 + x_2 \le 2,\\
          & x_1, x_2 \ge 0
\end{aligned}
\]
menghasilkan tableau berikut setelah satu pivot:
\begin{center}
\renewcommand{\arraystretch}{1.3}
\begin{tabular}{c|ZZZZ|Z}
\toprule
Basis & $x_1$ & $x_2$ & $s_1$ & $s_2$ & RHS\\
\midrule
$z$   & $0$ & $-2$ & $1$ & $0$ & $1$\\
\midrule
$x_1$ & $1$ & $-1$ & $1$ & $0$ & $1$\\
$s_2$ & $0$ & $0$  & $1$ & $1$ & $3$\\
\bottomrule
\end{tabular}
\end{center}
\begin{enumerate}
\item Temukan ketakterbatasan: identifikasi variabel masuk dan tunjukkan bahwa uji rasio tidak menemukan baris pembatas.
\item Sertifikasi: tetapkan variabel masuk sama dengan $t \ge 0$, tuliskan keluarga solusi $\x(t)$ yang dihasilkan dalam keempat variabel, dan verifikasi melalui substitusi langsung bahwa $\x(t)$ memenuhi kedua kendala asli dan syarat nonnegativitas untuk setiap $t \ge 0$.
\item Tunjukkan bahwa $z(t) \to \infty$ ketika $t \to \infty$, dan gambarkan daerah layak untuk menjelaskan secara geometris arah sinar sertifikat.
\end{enumerate}
\exrefs{\S\ref{sec:tableau-optimality}; Contoh~\ref{exa:simplex-unbounded}}
\end{ex}

\subsubsection*{Solusi Pilihan}

\begin{solution}(Latihan~\ref{ex:tableau-reading})
Variabel basis adalah label baris $x_1$ dan $x_2$; variabel nonbasis adalah
$s_1$ dan $s_2$. Menetapkan variabel nonbasis sama dengan nol dan membaca kolom
RHS memberikan solusi basis layak
\[
x_1 = 6, \quad x_2 = 4, \quad s_1 = s_2 = 0, \qquad z = 40.
\]
Entri baris objektif untuk variabel nonbasis adalah $3$ dan $1$, keduanya tak
negatif. Dalam konvensi ini, entri positif berarti bahwa menaikkan variabel
nonbasis yang bersesuaian akan \emph{menurunkan} $z$, sehingga tidak ada pivot
yang dapat memperbaiki objektif dan tableau tersebut optimal.
\end{solution}

\begin{solution}(Latihan~\ref{ex:tableau-ratio-test})
Entri baris objektif yang paling negatif adalah $-4$, sehingga $x_1$ masuk.
Uji rasio membagi setiap RHS dengan entri positif pada kolom $x_1$:
\[
\frac{6}{2} = 3, \qquad \frac{8}{4} = 2, \qquad \frac{12}{3} = 4.
\]
Rasio minimum $2$ terjadi pada baris $s_2$, sehingga $s_2$ keluar dan elemen
pivotnya adalah $4$ pada baris $s_2$ di kolom $x_1$.
\end{solution}

\begin{solution}(Latihan~\ref{ex:tableau-solve-lp})
Setelah menambahkan variabel slack $s_1, s_2$, tableau awal memiliki baris objektif $(-1, -2, 0, 0 \mid 0)$.

\textbf{Pivot pertama:} $y$ masuk ($-2$ paling negatif). Rasio: $4/1 = 4$ dan $6/3 = 2$, jadi $s_2$ keluar; elemen pivot adalah $3$ di baris $s_2$. Membagi baris pivot dengan $3$ dan menolkan entri lain pada kolom $y$ menghasilkan
\begin{center}
\renewcommand{\arraystretch}{1.3}
\begin{tabular}{c|ZZZZ|Z}
\toprule
Basis & $x$ & $y$ & $s_1$ & $s_2$ & RHS\\
\midrule
$z$   & $-\tfrac13$ & $0$ & $0$ & $\tfrac23$ & $4$\\
\midrule
$s_1$ & $\tfrac23$ & $0$ & $1$ & $-\tfrac13$ & $2$\\
$y$   & $\tfrac13$ & $1$ & $0$ & $\tfrac13$  & $2$\\
\bottomrule
\end{tabular}
\end{center}

\textbf{Pivot kedua:} $x$ masuk ($-\tfrac13$). Rasio: $2/\tfrac23 = 3$ dan $2/\tfrac13 = 6$, sehingga $s_1$ keluar; elemen pivot adalah $\tfrac23$ di baris $s_1$. Tableau yang dihasilkan adalah
\begin{center}
\renewcommand{\arraystretch}{1.3}
\begin{tabular}{c|ZZZZ|Z}
\toprule
Basis & $x$ & $y$ & $s_1$ & $s_2$ & RHS\\
\midrule
$z$   & $0$ & $0$ & $\tfrac12$ & $\tfrac12$ & $5$\\
\midrule
$x$ & $1$ & $0$ & $\tfrac32$ & $-\tfrac12$ & $3$\\
$y$ & $0$ & $1$ & $-\tfrac12$ & $\tfrac12$ & $1$\\
\bottomrule
\end{tabular}
\end{center}
Setiap entri baris objektif tak negatif, jadi tableau tersebut optimal: $(x, y) = (3, 1)$ dengan $z^* = 5$.
\end{solution}

\begin{solution}(Latihan~\ref{ex:tableau-two-pivot-solve})
Setelah menambahkan variabel slack $s_1, s_2$, tableau awal adalah
\begin{center}
\renewcommand{\arraystretch}{1.3}
\begin{tabular}{c|ZZZZ|Z}
\toprule
Basis & $x$ & $y$ & $s_1$ & $s_2$ & RHS\\
\midrule
$z$   & $-3$ & $-2$ & $0$ & $0$ & $0$\\
\midrule
$s_1$ & $2$ & $1$ & $1$ & $0$ & $10$\\
$s_2$ & $1$ & $3$ & $0$ & $1$ & $15$\\
\bottomrule
\end{tabular}
\end{center}

\textbf{Pivot pertama: } $x$ masuk ($-3$ paling negatif). Rasio: $10/2 = 5$ dan $15/1 = 15$, jadi $s_1$ keluar; elemen pivot adalah $2$ di baris $s_1$. Kalikan baris pivot dengan $\tfrac12$, lalu nolkan entri lain pada kolom $x$ untuk memperoleh
\begin{center}
\renewcommand{\arraystretch}{1.3}
\begin{tabular}{c|ZZZZ|Z}
\toprule
Basis & $x$ & $y$ & $s_1$ & $s_2$ & RHS\\
\midrule
$z$   & $0$ & $-\tfrac12$ & $\tfrac32$ & $0$ & $15$\\
\midrule
$x$   & $1$ & $\tfrac12$ & $\tfrac12$ & $0$ & $5$\\
$s_2$ & $0$ & $\tfrac52$ & $-\tfrac12$ & $1$ & $10$\\
\bottomrule
\end{tabular}
\end{center}

\textbf{Pivot kedua:} $y$ masuk ($-\tfrac12$). Rasio: $5/\tfrac12 = 10$ dan $10/\tfrac52 = 4$, sehingga $s_2$ keluar; elemen pivot adalah $\tfrac52$ di baris $s_2$. Tableau yang dihasilkan adalah
\begin{center}
\renewcommand{\arraystretch}{1.3}
\begin{tabular}{c|ZZZZ|Z}
\toprule
Basis & $x$ & $y$ & $s_1$ & $s_2$ & RHS\\
\midrule
$z$   & $0$ & $0$ & $\tfrac75$ & $\tfrac15$ & $17$\\
\midrule
$x$ & $1$ & $0$ & $\tfrac35$ & $-\tfrac15$ & $3$\\
$y$ & $0$ & $1$ & $-\tfrac15$ & $\tfrac25$ & $4$\\
\bottomrule
\end{tabular}
\end{center}
Setiap entri baris $z$ tak negatif, jadi tableau tersebut optimal: $(x, y) = (3, 4)$ dengan $z^* = 17$ dan $s_1 = s_2 = 0$ (kedua kendala aktif pada optimum).
\end{solution}
